\documentclass[11pt]{article}
\usepackage[T1]{fontenc}
\usepackage{lmodern}
\usepackage{amsmath,amssymb,amsthm}
\usepackage[margin=1in]{geometry}
\usepackage{microtype}
\usepackage[colorlinks=true,linkcolor=blue,citecolor=blue,urlcolor=blue]{hyperref}
\hypersetup{
 pdftitle={Equidistribution of Prime-Degree Torus Packets with Arbitrary Local Type},
 pdfauthor={OpenAI},
 pdfsubject={Full-packet equidistribution for arbitrary local module types},
 pdfkeywords={torus packets, equidistribution, prime degree, nonmaximal orders, entropy}
}
\newtheorem{theorem}{Theorem}[section]
\newtheorem{proposition}[theorem]{Proposition}
\newtheorem{lemma}[theorem]{Lemma}
\newtheorem{corollary}[theorem]{Corollary}
\theoremstyle{remark}
\newtheorem{remark}[theorem]{Remark}
\newcommand{\R}{\mathbb R}
\newcommand{\Z}{\mathbb Z}
\newcommand{\Q}{\mathbb Q}
\DeclareMathOperator{\N}{N}
\DeclareMathOperator{\vol}{vol}
\DeclareMathOperator{\Disc}{Disc}
\DeclareMathOperator{\SL}{SL}
\DeclareMathOperator{\diag}{diag}
\newcommand{\kappaK}{\kappa_K}
\newcommand{\dd}{\,d}
\newcommand{\one}{\mathbf 1}
\newcommand{\eps}{\varepsilon}
\newcommand{\cO}{\mathcal O}
\newcommand{\cV}{\mathcal V}
\numberwithin{equation}{section}
\title{Equidistribution of Prime-Degree Torus Packets\\
with Arbitrary Local Type}
\author{OpenAI}
\date{September 24, 2026}
\begin{document}
\maketitle
\begin{abstract}
We prove the packet form of the higher-dimensional Duke equidistribution
problem for totally real fields of any fixed prime degree at least five,
allowing arbitrary local homothety types of full lattices. As the
multiplier-order discriminant tends to infinity, the volume-weighted packet
measures converge to Haar probability measure, with no escape of mass.
\end{abstract}

\section{Introduction}\label{sec:intro}

Compact diagonal orbits in the space of unimodular lattices encode
ideal classes and units of totally real number fields. The higher-dimensional
Duke equidistribution problem asks whether arithmetically complete families
of these orbits become uniformly distributed as their discriminants grow.
The grouping matters: one local homothety type specifies a complete packet,
whereas a multiplier order alone can combine several packets. We prove the
prime-degree packet statement for every such local type.

Fix a prime $n\ge5$, and write
\[
 X_n=\SL_n(\Z)\backslash\SL_n(\R),\qquad
 A_n=\{\diag(e^{u_1},\ldots,e^{u_n}):\textstyle\sum_j u_j=0\}.
\]
We regard $X_n$ as the space of covolume-one row lattices, with the
right action of $A_n$, and let $m_n$ be its $\SL_n(\R)$-invariant
probability measure. Haar measure on $A_n$ is
$du_1\cdots du_{n-1}$ in the displayed logarithmic coordinates.

Let $K$ be a totally real field of degree $n$, let
$\sigma=(\sigma_1,\ldots,\sigma_n)$ order its real embeddings, and let
$M\subset K$ be a free rank-$n$ $\Z$-submodule spanning $K$ over $\Q$.
We call such an $M$ a full lattice. Its multiplier order and discriminant
are
\[
 \cO(M)=\{\alpha\in K:\alpha M\subset M\},\qquad
 D(M)=|\Disc(K)|[\cO_K:\cO(M)]^2.
\]
Two such lattices are \emph{locally homothetic} if
$M'_\ell=c_\ell M_\ell$ for every rational prime $\ell$, where
$c_\ell\in(K\otimes_\Q\Q_\ell)^\times$ and subscripts denote completion.
Choose representatives $M'$ for their $K^\times$-homothety classes
within the local type of $M$, and normalize
\[
 \Lambda_{M',\sigma}
 =\operatorname{covol}(\sigma(M'))^{-1/n}\sigma(M').
\]
The packet $\mathcal P_{K,M,\sigma}$ consists of the distinct orbits
$\Lambda_{M',\sigma}wA_n$, with
\[
 w\in W_n=\{\diag(\epsilon_1,\ldots,\epsilon_n):
       \epsilon_j\in\{\pm1\},\ \textstyle\prod_j\epsilon_j=1\}.
\]
These orbits are compact and their number is finite; the
parametrization in Section~\ref{sec:packets} also proves these facts.
For an orbit $O$, let $\nu_O$ be its invariant probability and
$\vol(O)$ its volume for the specified Haar measure. Set
\begin{equation}\label{eq:packet-probability}
 \mu_{K,M,\sigma}
 =\frac{\sum_{O\in\mathcal P_{K,M,\sigma}}\vol(O)\nu_O}
        {\sum_{O\in\mathcal P_{K,M,\sigma}}\vol(O)}.
\end{equation}

\begin{theorem}\label{thm:main}
For every fixed prime $n\ge5$ and every sequence
$(K_i,M_i,\sigma_i)$ as above with \mbox{$D(M_i)\to\infty$},
\[
 \mu_{K_i,M_i,\sigma_i}\longrightarrow m_n.
\]
The convergence is weak convergence of probability measures. In
particular, it holds against every continuous compactly supported
function, and no mass escapes into the cusp.
\end{theorem}

Theorem~\ref{thm:main} permits nonmaximal orders of unbounded index
and arbitrary local module types. The field discriminants may remain
bounded or grow independently of those indices. The statement is
separate for each local homothety type; it makes no assertion about an
individual orbit, a convergence rate, or composite degrees.

\subsection*{The maximal-order specialization}

Take $M=\cO_K$. Every lattice in its local type is a fractional
$\cO_K$-ideal, and its homothety class is an ordinary ideal class.
Thus the packet is obtained by embedding one representative of each
ideal class, normalizing its covolume to one, and taking the distinct
sign translates of its positive diagonal orbit. For every sequence of
totally real fields of degree $n$ with $|\Disc(K)|\to\infty$, the
volume-weighted measures of these packets converge to $m_n$ by
Theorem~\ref{thm:main}. This construction also explains why an individual
ideal class is not the object being averaged.

For general $M$, the same maximal-order ideal classes remain the global
parameters, but each carries a finite set of local module choices.
Proposition~\ref{prop:adelic-packet} in Appendix~\ref{subsec:adelic-packet}
identifies this construction with the
pushforward of adelic torus Haar probability, independently of the
rational basis and finite-lattice presentation. The theorem therefore
includes all these presentations of the maximal-order specialization,
as well as arbitrary noninvertible modules over nonmaximal orders.

\subsection*{Context and prior work}

Linnik's ergodic approach connected arithmetic distribution problems
with dynamics \cite{Linnik1968}. Duke proved equidistribution of the
length-weighted closed geodesics associated with fundamental
real-quadratic discriminants on the modular surface
\cite[Theorem~1]{Duke1988}. In higher rank, Einsiedler, Lindenstrauss,
Michel, and Venkatesh developed the relationship between arithmetic
separation and entropy for periodic torus orbits \cite{ELMV2009}.
Their cubic equidistribution theorem combines subconvexity, local
analysis, and measure classification
\cite[Theorem~1.4]{ELMV2011}. Their packet construction already
distinguishes local homothety from the coarser condition of having the
same multiplier order \cite[Section~5.5]{ELMV2011}.

The higher-prime-degree extension described in
\cite[Section~1.6.3]{ELMV2011} is conditional on subconvexity.
Khayutin subsequently obtained entropy bounds using arithmetic
separation and Galois invariants, under maximal-order and
two-transitivity hypotheses \cite[Theorem~1.1]{Khayutin2019}.
Lemke Oliver, Thorner, and Zaman proved equidistribution for sequences
of orders in totally real fields of each fixed prime degree at least five,
outside a small
exceptional family of fields \cite[Theorem~2.3]{LOTZ2024}.
Their result allows nonmaximal orders; the exceptional family has at
most $O_{n,\eps}(X^\eps)$ fields of discriminant at most $X$.
Theorem~\ref{thm:main} treats every field and every full-lattice local
type in the stated prime degrees. It gives an unconditional affirmative
answer to this packet form of the higher-dimensional Duke problem,
including sequences in a fixed field with increasing order index.

\subsection*{The arithmetic ingredient}

The proof follows the counting-to-rigidity strategy of
Einsiedler--Lindenstrauss--Michel--Venkatesh
\cite[Section~2.7.1]{ELMV2011}: packet averages of lattice-vector
counts control the cusp and small ordinary balls, and these bounds
feed into positive-entropy rigidity. Their cubic argument obtains
asymptotics for these averages using global and local subconvexity
\cite[Sections~2.7.2--2.7.3]{ELMV2011}. For the same rigidity step,
we obtain the needed upper bounds by retaining the Dedekind-zeta
residue in the unfolded ideal count. Stark's exceptional-zero descent
\cite{Stark1974}, together with Shiu's theorem in Pollack's uniform
formulation \cite{Shiu1980,Pollack2020}, supplies the required bounds
without a subconvexity hypothesis. Odd degree excludes the quadratic
subfields that occur in Stark's exceptional-zero alternative.

For an arbitrary prescribed local type, we keep maximal-order ideal
counting but insert local containment probabilities: at each prime,
the weight records the probability that a vector belongs to a randomly
chosen unit translate of the prescribed module.
The essential estimate is a small positive moment of the local norm,
uniform over full local modules whose maximal-order span is the full
ring. Its proof uses a unit in each component projection, not
independence of the components or invertibility of the module.
The unit-mass argument in \cite[Lemma~9.7]{ELMV2011} is an antecedent;
here we prove the exact shell normalization and the positive moment
for arbitrary local modules with this full-span condition.

The exact unfolding belongs to the classical Hecke method, presented
in the packet setting in \cite[Section~10]{ELMV2011}. The dynamical
conclusion uses the two-sided tube criterion of
\cite[Corollary~3.3]{ELMV2009} and the prime-degree measure
classification of Einsiedler, Katok, and Lindenstrauss
\cite[Corollary~1.4]{EKL2006}. We give the residue estimate,
local-module bounds, and the passage to entropy in almost every
ergodic component below.

\subsection*{Proof outline}

Put $R=\cO_K$, $D=|\Disc(K)|$, and
$q=[RM:M]$. The natural counting scale is $Q=q\sqrt D$, and
$D(M)\le Q^{2n}$. Section~\ref{sec:packets} expresses the packet
mean of a lattice-vector count as a maximal-order ideal sum with
local weights. Section~\ref{sec:local} proves an exact total-mass
identity and a power-saving tail for these weights.
Section~\ref{sec:analytic} combines them with a residue-relative
zeta estimate and Shiu's theorem. Removing the Euler factors at
primes dividing $q$ produces precisely the factor needed to cancel
the local normalization.

The resulting count controls both short lattice vectors and small
boxes about vectors with nonzero coordinates. Section~\ref{sec:geometry}
deduces tightness and an $O(r^n)$ bound for the mass of ordinary balls
in any limit measure. Since $A_n$ has dimension $n-1$, this bound
forces positive entropy in almost every $A_n$-ergodic component.
Section~\ref{sec:entropy} gives this implication and applies the
prime-degree measure-classification theorem. In that step it is
essential to prove positive entropy for every invariant probability
dominated by a finite multiple of the limit, not only for the limit
itself: this rules out a positive-weight family of zero-entropy
components.

\section{The packet as a weighted ideal sum}\label{sec:packets}

We first identify the packet measure exactly, including its signs and
orbit volumes. Throughout the proof, $K$ is totally real of the fixed
prime degree $n$, and
\[
 R=\cO_K,\quad E=R^\times,\quad D=|\Disc(K)|,\quad
 \kappaK=\operatorname*{Res}_{s=1}\zeta_K(s).
\]
The absolute norm $\N I$ of a nonzero fractional $R$-ideal is extended
multiplicatively from integral ideals. Thus
$\operatorname{covol}(\sigma(I))=\sqrt D\,\N I$.
We use standard ideal theory and the unit theorem
\cite[Chapters~3--5]{Milne2020}, together with the analytic
class-number formula \cite[Theorem~19.12]{Sutherland2021}.
The classical packet interpretation is discussed in
\cite[Sections~2 and~5]{ELMV2011}.

\subsection{Saturation and local choices}

Let $J_0=RM$, $q=[J_0:M]$, and $Q=q\sqrt D$. For a rational prime
$\ell$, write $(J_0)_\ell=j_\ell R_\ell$, and put
$L_\ell=j_\ell^{-1}M_\ell$. Then
\begin{equation}\label{eq:local-normalization}
 L_\ell\subset R_\ell,\qquad R_\ell L_\ell=R_\ell,\qquad
 [R_\ell:L_\ell]=q_\ell,\qquad q=\prod_\ell q_\ell.
\end{equation}
Set $S=\{\ell:\ell\mid q\}$; off $S$, $L_\ell=R_\ell$.
Since $qJ_0\subset M$, one has $qR\subset\cO(M)$.
Conversely, a multiplier of $M$ has an integral endomorphism matrix
on $M$ and hence is an algebraic integer. It follows that
\begin{equation}\label{eq:discriminant-scale}
 qR\subset\cO(M)\subset R,\qquad
 D(M)\le Dq^{2n}\le Q^{2n}.
\end{equation}
Only the last implication, $D(M)\to\infty\Rightarrow Q\to\infty$,
will be needed.

These quantities are constant within the prescribed local type.
Indeed, if $M'_\ell=c_\ell M_\ell$, multiplication by $c_\ell$
identifies $(RM)_\ell/M_\ell$ with $(RM')_\ell/M'_\ell$.
It also identifies the local multiplier orders, since multiplication
in $K\otimes_\Q\Q_\ell$ is commutative. The multiplier order commutes
with completion: in a $\Z$-basis of $M$, its defining condition is
that the matrix of multiplication have integral entries, and the
same condition over $\Z_\ell$ defines the local order. Thus $q$,
$Q$, $\cO(M)$, and $D(M)$ are unchanged by local homothety.
Different local types can nevertheless have the same multiplier order.

Choose one fractional ideal $I$ in each ordinary ideal class of $R$.
Let $\mathcal N(I)$ be the set of lattices $N$ locally homothetic to $M$
with $RN=I$. If $I_\ell=i_\ell R_\ell$, its possible completions are
exactly
\begin{equation}\label{eq:local-choices}
 N_\ell=i_\ell z_\ell L_\ell,\qquad z_\ell\in R_\ell^\times.
\end{equation}
Indeed, taking maximal-order spans in a local homothety forces the
remaining multiplier to be a unit. Conversely, these local lattices
lie between $qI_\ell$ and $I_\ell$ and patch uniquely by the primary
decomposition of $I/qI$. Thus $\mathcal N(I)$ is finite,
$[I:N]=q$ for all its members, and its cardinality
$T=|\mathcal N(I)|$ is independent of $I$. The label count $T$ need
not equal the saturation index $q$. Local unit Haar probability induces
uniform probability on each finite set of local choices, so the choices at different
primes are independent. Every global homothety class in the
packet occurs in these sets, by moving its maximal-order span
to the chosen ideal-class representative.

\subsection{Real parameters and measure}

Let
\[
 C=\{t\in(\R^\times)^n:\textstyle\prod_j|t_j|=1\}.
\]
Give $C$ counting measure on its $2^n$ sign components and measure
$du_1\cdots du_{n-1}$ on the coordinates $u_j=\log|t_j|$.
The embedding of $E$ is a discrete cocompact subgroup of $C$.
If $F$ is a measurable fundamental domain, the unit theorem and
the totally real case of the analytic class-number formula
\cite[Theorem~19.12]{Sutherland2021} give
\begin{equation}\label{eq:unit-volume}
 \vol(F)=2^{n-1}R_K,\qquad
 h_K\vol(F)=\sqrt D\,\kappaK,
\end{equation}
where $h_K$ is the ordinary class number and $R_K$ the regulator
in deleted logarithmic coordinates. The factor $2^{n-1}$ comes
from the $2^n$ signs and the logarithmic kernel $\{1,-1\}$.
In the class-number formula this is the factor $2^n/w_K$, since
the roots of unity in a totally real field are exactly $\{1,-1\}$.

\begin{lemma}\label{lem:packet-model}
The maps
\[
 (N,t)\longmapsto (Q\N I)^{-1/n}\sigma(N)t
 \qquad (N\in\mathcal N(I),\ t\in C)
\]
induce a bijection from the disjoint union of the quotients by
\[
 \epsilon:(N,t)\longmapsto
       (\epsilon N,\sigma(\epsilon)^{-1}t),\qquad \epsilon\in E,
\]
onto the union of the packet orbits.
Counting measure on the $N$ labels times the stated measure on $C$
descends to the sum of the orbit Haar measures. In particular,
the packet probability is obtained by averaging uniformly over $I$
and $\mathcal N(I)$ and integrating over $F$ with measure
$dt/\vol(F)$.
\end{lemma}

\begin{proof}
Surjectivity follows from \eqref{eq:local-choices}. A global
homothety changes a normalized embedded lattice by an element
of $C$. Since $n$ is odd, each pair of opposite sign patterns has
exactly one representative in $W_n$; a common negative sign
does not change a lattice.

For injectivity, suppose two displayed normalized lattices coincide.
Absorb the ratio of the normalizing scalars into the real diagonal
map carrying $\sigma(N)$ onto $\sigma(N')$. Applying this map
to one nonzero $\alpha\in N$ gives $\sigma(\beta)$ for some
$\beta\in N'$. None of the coordinates of $\sigma(\alpha)$
vanishes, so the map is multiplication by $\sigma(\beta/\alpha)$.
Thus $N'=bN$ for $b\in K^\times$. Their maximal-order spans
represent the same ideal class and therefore equal the same
chosen $I$. Hence $bI=I$, so $b\in E$, and the parameters differ
by exactly the stated action.

Combine an $N$ label with a sign component of $C$. Its stabilizer in $E$
is the totally positive subgroup of
$\{\epsilon\in E:\epsilon N=N\}$, and has finite index in $E$.
It acts faithfully by logarithmic translations. By injectivity,
these are all periods of the corresponding $A_n$-orbit.
Each orbit of labels therefore contributes exactly one compact
$A_n$-orbit, with precisely its prescribed Haar volume.
This proves both finiteness and the measure assertion.
One may equally use $\mathcal N(I)\times F$ as a fundamental
domain, since the $E$-action preserves counting measure times $dt$.
The total measure of these fundamental domains is
$h_KT\vol(F)$. Dividing by this quantity gives precisely the
uniform averages over ideal classes and finite labels, followed
by integration against $dt/\vol(F)$.
\end{proof}

\subsection{Containment probabilities and unfolding}

Write $\mathfrak p\mid\ell$ for the primes of $R$ over $\ell$,
with ramification indices $e_{\mathfrak p}$ and residue degrees
$f_{\mathfrak p}$. Normalize $v_{\mathfrak p}$ by a uniformizer,
so $v_{\mathfrak p}(\ell)=e_{\mathfrak p}$ and
$\sum_{\mathfrak p\mid\ell}e_{\mathfrak p}f_{\mathfrak p}=n$.
For a tuple $\mathbf k_\ell=(k_{\mathfrak p})_{\mathfrak p\mid\ell}$
of nonnegative integers, define
\[
 P_\ell(\mathbf k_\ell)
 =\Pr_{z\in R_\ell^\times}\{x\in zL_\ell\},
\]
where $x$ has those component valuations and the probability is
normalized unit Haar measure. The definition is independent of $x$
because the units act transitively on each valuation shell.
For an integral nonzero $R$-ideal $\mathfrak a$, set
\[
 P(\mathfrak a)=\prod_{\ell\in S}
       P_\ell((v_{\mathfrak p}(\mathfrak a))_{\mathfrak p\mid\ell}).
\]
For $\alpha\in I\setminus\{0\}$, the probability that a uniform
$N\in\mathcal N(I)$ contains $\alpha$ is
$P((\alpha)I^{-1})$.

For a nonnegative Borel function $f$ on $\R^n$, put
$E_f(\Lambda)=\sum_{0\ne v\in\Lambda}f(v)$.
The product-fiber integral below integrates $f$ over the $C$-orbit
with absolute coordinate product $z$, using logarithmic Haar measure
and summing over all sign components:
\begin{equation}\label{eq:product-fiber}
 \cV_f(z)=\sum_{\omega\in\{\pm1\}^n}
 \int_{\R^{n-1}}
 f\left(\omega_1e^{u_1},\ldots,\omega_{n-1}e^{u_{n-1}},
      \omega_n z e^{-\sum_{j<n}u_j}\right)
 \dd u_1\cdots\dd u_{n-1},\qquad z>0.
\end{equation}

\begin{proposition}\label{prop:unfolding}
Writing $\mu_*=\mu_{K,M,\sigma}$, one has
\begin{equation}\label{eq:unfolding}
 \int_{X_n}E_f\dd\mu_*
 =\frac1{\sqrt D\,\kappaK}
       \sum_{\mathfrak a\subset R,\ \mathfrak a\ne0}
           P(\mathfrak a)\cV_f(\N\mathfrak a/Q).
\end{equation}
The identity holds for nonnegative extended integrals.
\end{proposition}

\begin{proof}
Lemma~\ref{lem:packet-model} expresses the left side as
\[
 \frac1{h_K\vol(F)}
 \sum_I\int_F\sum_{\alpha\in I\setminus\{0\}}
   P((\alpha)I^{-1})
   f((Q\N I)^{-1/n}\sigma(\alpha)t)\dd t.
\]
The map $\alpha\mapsto(\alpha)I^{-1}$ runs over the integral
ideals in the inverse class of $I$, each fiber being one free
$E$-orbit. Summing that orbit and integrating over $F$ unfolds
to an integral over $C$, by Tonelli. The absolute coordinate
product of the normalized vector is
$\N((\alpha)I^{-1})/Q$. Translating the first $n-1$ logarithms
and relabeling the signs identifies the $C$-integral with
\eqref{eq:product-fiber}, with Jacobian one.
Every integral ideal occurs in exactly one inverse class.
Equation~\eqref{eq:unit-volume} supplies the coefficient in
\eqref{eq:unfolding}.
\end{proof}

When $M$ is an $R$-ideal, $q=1$ and $P(\mathfrak a)=1$.
Thus \eqref{eq:unfolding} recovers the maximal-order formula.

\section{Uniform control of the local weights}\label{sec:local}

The ideal sum in Proposition~\ref{prop:unfolding} involves arbitrary
local modules. We now isolate two estimates that are uniform in
those modules and in their indices. The first gives the total
normalization; the second makes large local norm contributions
negligible.

Let $\mathbf k=(\mathbf k_\ell)_{\ell\in S}$ range over nonnegative
valuation tuples, and write
\[
 d_\ell(\mathbf k_\ell)
   =\ell^{\sum_{\mathfrak p\mid\ell}f_{\mathfrak p}k_{\mathfrak p}},
 \qquad d(\mathbf k)=\prod_{\ell\in S}d_\ell(\mathbf k_\ell),
 \qquad P(\mathbf k)=\prod_{\ell\in S}P_\ell(\mathbf k_\ell).
\]
Thus $d$ is the norm of the $S$-supported part of an ideal. Also set
\[
 Z_\ell=\prod_{\mathfrak p\mid\ell}
          (1-\ell^{-f_{\mathfrak p}})^{-1},\qquad
 Z_S=\prod_{\ell\in S}Z_\ell.
\]
Empty products are one.

\begin{lemma}\label{lem:local-moment}
For the weights just defined,
\begin{equation}\label{eq:shell-sum}
 \sum_{\mathbf k}\frac{P(\mathbf k)}{d(\mathbf k)}=\frac{Z_S}{q}.
\end{equation}
If $s=1/(2n)$, then, for every $\eps>0$,
\begin{equation}\label{eq:weight-moment}
 \sum_{\mathbf k}P(\mathbf k)d(\mathbf k)^{s-1}
       \ll_{n,\eps}\frac{q^\eps}{q}.
\end{equation}
In particular,
\begin{equation}\label{eq:weight-tail}
 \sum_{d(\mathbf k)>Q^{1/8}}\frac{P(\mathbf k)}{d(\mathbf k)}
       \ll_n \frac1q Q^{-s/16}.
\end{equation}
\end{lemma}

\begin{proof}
Normalize additive Haar measure on $R_\ell$ to have mass one.
The valuation shell $\mathbf k_\ell$ has mass
$Z_\ell^{-1}/d_\ell(\mathbf k_\ell)$. Its conditional law is
invariant under the transitive unit action, so $P_\ell$ is
exactly the fraction of this shell lying in $L_\ell$.
Since $L_\ell$ has mass $1/q_\ell$, summing over shells and
then multiplying over $S$ proves \eqref{eq:shell-sum}.
Elements with a zero component have additive measure zero.

Choose an additive-Haar uniform element of $L_\ell$, and denote
its component valuations by $v_{\mathfrak p}$. The condition
$R_\ell L_\ell=R_\ell$ implies that the projection to every
component contains a unit. Modulo the $k$th power of its maximal
ideal, the projection therefore contains all $\Z_\ell$ multiples
of that unit, at least $\ell^{\lceil k/e_{\mathfrak p}\rceil}$
elements. Uniform measure on $L_\ell$ induces uniform measure
on this finite additive image. Consequently
\begin{equation}\label{eq:valuation-tail}
 \Pr(v_{\mathfrak p}\ge k)
       \le\ell^{-\lceil k/e_{\mathfrak p}\rceil}\qquad(k\ge1).
\end{equation}
This argument does not require the component projections to be
independent, or even to equal their full integer rings.

Put $j=\max_{\mathfrak p\mid\ell}
\lceil v_{\mathfrak p}/e_{\mathfrak p}\rceil$. Then
$d_\ell(\mathbf v_\ell)\le\ell^{nj}$. If $j\ge a\ge1$, some
$v_{\mathfrak p}$ is at least $e_{\mathfrak p}(a-1)+1$.
Equation~\eqref{eq:valuation-tail} and a union bound give
$\Pr(j\ge a)\le n\ell^{-a}$. Hence
\begin{equation}\label{eq:local-moment}
 \mathbb E_{L_\ell}d_\ell(\mathbf v_\ell)^s
 \le1+n\sum_{a\ge1}\ell^{(ns-1)a}
 \le1+C_n\ell^{-1/2}.
\end{equation}
Using the shell masses once more yields the exact factorization
\[
 \sum_{\mathbf k}P(\mathbf k)d(\mathbf k)^{s-1}
   =\frac1q\prod_{\ell\in S}
       \left(Z_\ell\,\mathbb E_{L_\ell}d_\ell(\mathbf v_\ell)^s\right).
\]
Each factor is at most
$(1-\ell^{-1})^{-n}(1+C_n\ell^{-1/2})$. For fixed $\eps>0$
this is at most $\ell^\eps$ at all sufficiently large primes;
the finitely many remaining primes contribute a constant
depending only on $n,\eps$. As $\prod_{\ell\in S}\ell\le q$,
this proves \eqref{eq:weight-moment}. Finally, take
$\eps=s/16$ and use $q\le Q$ to obtain
\[
 \sum_{d>Q^{1/8}}\frac{P(\mathbf k)}d
 \le Q^{-s/8}\sum_{\mathbf k}P(\mathbf k)d^{s-1}
 \ll_n q^{-1}Q^{-s/16}.
 \qedhere
\]
\end{proof}

\section{Residue-relative weighted ideal counts}\label{sec:analytic}

Write
\[
 a_K(m)=\#\{\mathfrak a\subseteq R:\N\mathfrak a=m\},
 \qquad \kappaK=\operatorname*{Res}_{s=1}\zeta_K(s)>0.
\]
We first retain the residue in the analytic estimates. All constants
subscripted by $n$ are uniform over the fields under consideration.

\begin{lemma}\label{lem:residue}
For $D$ sufficiently large in terms of $n$,
\begin{equation}\label{eq:residue}
 h\zeta_K(1+h)\asymp_n\kappaK
 \quad\left(0<h\le\frac4{\log D}\right),
 \qquad \kappaK^{-1}\ll_n\log D.
\end{equation}
For the remaining fields, the first comparison holds uniformly for
sufficiently small positive $h$, and $\kappaK^{-1}\ll_n1$.
\end{lemma}

\begin{proof}
Put $L=\log D$. Stark's Lemma~3 gives at most one zero in
$\operatorname{Re}s\ge1-(4L)^{-1}$,
$|\operatorname{Im}s|\le(4L)^{-1}$; such a zero is real and simple.
His Lemma~8, which also applies to nonnormal fields, says that a real
zero $\beta\ge1-(4n!L)^{-1}$ forces a quadratic subfield of $K$
\cite[Lemmas~3 and~8]{Stark1974}. Since $n$ is odd, there is no such
subfield. Consequently, for a constant $c_n>0$, every nontrivial zero
$\rho$ satisfies $|1-\rho|\ge c_n/L$.

The real logarithmic derivative of the completed zeta function gives
\[
 S_K(\sigma):=\sum_\rho\operatorname{Re}\frac1{\sigma-\rho}
 =\frac1\sigma+\frac1{\sigma-1}+\frac L2
   +\frac{\zeta_K'}{\zeta_K}(\sigma)+G_n(\sigma),\qquad \sigma>1,
\]
where zeros are counted with multiplicity and
\[
 G_n(\sigma)=\frac n2
 \left(\frac{\Gamma'(\sigma/2)}{\Gamma(\sigma/2)}-\log\pi\right).
\]
This is the Hadamard identity in
\cite[Equation~(9)]{Stark1974}. The zero terms are nonnegative,
$G_n$ is bounded on $[1,2]$, and the Euler product gives
$\zeta_K'/\zeta_K\le0$ there. Thus, at $\sigma_0=1+L^{-1}$,
$S_K(\sigma_0)\ll_n L$.

This bound persists uniformly for $1<\sigma\le1+4/L$.
Indeed, for $\rho=\beta+i\gamma$, set
$X=L(1-\beta)$, $Y=L\gamma$, and $U=L(\sigma-1)$.
Then $X\ge0$, $0<U\le4$, and $|X+iY|\ge c_n$. Hence
\[
 \frac{\operatorname{Re}(\sigma-\rho)^{-1}}
      {\operatorname{Re}(\sigma_0-\rho)^{-1}}
 =\frac{X+U}{X+1}
   \frac{(X+1)^2+Y^2}{(X+U)^2+Y^2}
 \le4(1+3/c_n)^2.
\]
Summing gives $0\le S_K(\sigma)\ll_n L$.
For $F_K(\sigma)=(\sigma-1)\zeta_K(\sigma)$, it follows that
\[
 \frac{F_K'}{F_K}(\sigma)
 =S_K(\sigma)-\frac1\sigma-\frac L2-G_n(\sigma)
 =O_n(L).
\]
Since $F_K$ extends analytically to $1$ with $F_K(1)=\kappaK>0$,
integration over an interval of length at most $4/L$ proves the
comparison in \eqref{eq:residue}. At $\sigma_0$, the inequality
$\zeta_K(\sigma_0)\ge1$ also gives $\kappaK\gg_n L^{-1}$.

For completeness, only finitely many fields of this fixed degree
have discriminant below the cutoff. One can see this directly here:
nonzero vectors of $D^{-1/(2n)}\sigma(R)$ have absolute coordinate
product at least $D^{-1/2}$, so Mahler's criterion puts these lattices
in a fixed compact set. Local basis lifts and a finite cover give
bases of $R$ with uniformly bounded conjugates. A basis element
outside $\Q$ generates $K$, since $n$ is prime; its monic integral
minimal polynomial therefore has bounded coefficients. There are
only finitely many such polynomials. The simple pole of each of
these finitely many zeta functions now gives the remaining claims
uniformly for sufficiently small $h>0$.
\end{proof}

We also use the following uniform form of Shiu's theorem
\cite{Shiu1980,Pollack2020}: if $b$ is nonnegative and multiplicative
and $b(m)\le d_n(m)$, where $d_n$ is the $n$-fold divisor function,
then
\begin{equation}\label{eq:shiu}
 \sum_{u-v<m\le u}b(m)
 \ll_n\frac v{\log u}
       \exp\left(\sum_{\ell\le u}\frac{b(\ell)}\ell\right)
 \qquad(u^{1/2}\le v\le u).
\end{equation}
Here $\ell$ denotes a rational prime, and the threshold for $u$
depends only on $n$. To check uniformity, the growth conditions in
\cite[Theorem~1.1]{Pollack2020}, with all excluded residue sets
empty, are supplied by
$d_n(\ell^j)=\binom{j+n-1}{n-1}\le n^j$ and
$d_n(m)\ll_{n,\eps}m^\eps$. Its strict interval condition
$u^\beta<v\le u$ is satisfied by taking $\beta=1/4$; thus
\eqref{eq:shiu} includes both endpoints of the displayed range.

\begin{proposition}\label{prop:weighted-count}
Fix $A_0\ge1$ and $0<\eta\le1$. For $Q$ sufficiently large in terms
of $n,A_0,\eta$,
\begin{equation}\label{eq:weighted-count}
 \sum_{x-y<\N\mathfrak a\le x}P(\mathfrak a)
 \ll_n \frac{\kappaK y}{q}
 \qquad
 (Q^{3/4}\le x\le A_0Q,\quad \eta x\le y\le x).
\end{equation}
The implied constant is independent of $A_0$ and $\eta$.
\end{proposition}

\begin{proof}
Set $b(m)=a_K(m)\one_{(m,q)=1}$. This function is nonnegative,
multiplicative, and bounded by $d_n$. Indeed, at every prime $\ell$,
including the ramified primes, the local generating series for
$a_K$ is $\prod_{\mathfrak p\mid\ell}(1-T^{f_{\mathfrak p}})^{-1}$;
it is coefficientwise bounded by $(1-T)^{-n}$.
Decomposing an ideal into its part over $S$ and its coprime part,
the sum in \eqref{eq:weighted-count} becomes
\begin{equation}\label{eq:weighted-split}
 \sum_{\mathbf k}P(\mathbf k)
       \sum_{(x-y)/d(\mathbf k)<m\le x/d(\mathbf k)}b(m).
\end{equation}

First suppose $d=d(\mathbf k)\le Q^{1/8}$, and put $u=x/d$.
Then $u\ge Q^{5/8}\ge D^{1/4}$. For $h=1/\log u$,
Lemma~\ref{lem:residue} gives
$\zeta_K(1+h)\ll_n\kappaK\log u$.
There is no upper restriction on $u$: increasing $u$ only decreases
$h$, so the lemma applies even when $D$ stays bounded and $q$ grows.
Positivity of the Euler logarithm yields
\[
 \sum_{\ell\le u}a_K(\ell)\ell^{-1-h}\le\log\zeta_K(1+h).
\]
Since $a_K(\ell)\le n$ and
$\sum_{\ell\le u}(\log\ell)/\ell\ll\log u$, replacing
$\ell^{-1-h}$ by $\ell^{-1}$ costs only $O_n(1)$.
The latter prime-sum bound follows, for example, by comparing
$\sum_{\ell\le u}\lfloor u/\ell\rfloor\log\ell$ with
$\log(\lfloor u\rfloor!)$. Therefore
\[
 \exp\left(\sum_{\ell\le u}\frac{a_K(\ell)}\ell\right)
 \ll_n\kappaK\log u.
\]
Furthermore,
\[
 \log Z_S=\sum_{\ell\in S}\frac{a_K(\ell)}\ell+O_n(1),
 \qquad
 \sum_{\substack{\ell\in S\\\ell>u}}
       \frac{a_K(\ell)}\ell
 \ll_n\frac{\log(2q)}{Q^{5/8}}\ll_n1.
\]
The first error is bounded by a constant times
$\sum_\ell\ell^{-2}$. Removing the primes of $S$ thus gives
\begin{equation}\label{eq:coprime-euler}
 \exp\left(\sum_{\ell\le u}\frac{b(\ell)}\ell\right)
 \ll_n\frac{\kappaK\log u}{Z_S}.
\end{equation}
For $Q$ sufficiently large at fixed $\eta$, the length
$v=y/d\ge\eta u$ is at least $u^{1/2}$.
Applying \eqref{eq:shiu}, then \eqref{eq:coprime-euler} and
\eqref{eq:shell-sum}, bounds the contribution of these $d$ by
\[
 C_n\frac{\kappaK y}{Z_S}
       \sum_{d(\mathbf k)\le Q^{1/8}}\frac{P(\mathbf k)}{d(\mathbf k)}
 \le C_n\frac{\kappaK y}{q}.
\]

For $d>Q^{1/8}$, terms with $d>x$ vanish because the inner sum
in \eqref{eq:weighted-split} then contains no positive integer.
For the remaining terms, $x/d\ge1$, and the elementary bound
$\sum_{m\le z}d_n(m)\le z(1+\log z)^{n-1}$ for $z\ge1$
follows by summing reciprocals in the first $n-1$ factors.
Using \eqref{eq:weight-tail}, with $s=1/(2n)$, the remaining
contribution to \eqref{eq:weighted-split} is at most
\[
 C_n\frac{x}{q}Q^{-s/16}(1+\log x)^{n-1}
 =o\left(\frac{\kappaK x}{q}\right).
\]
This is uniform in the stated range: Lemma~\ref{lem:residue}
gives $\kappaK^{-1}\ll_n\log(3D)$, and $D\le Q^2$.
At fixed $A_0,\eta$, increase the threshold until this term is at
most $\kappaK\eta x/q\le\kappaK y/q$. The implied constant has
not changed with $\eta$, proving the proposition.
\end{proof}

\section{Nonescape and ordinary ball bounds}\label{sec:geometry}

We apply the weighted count in two ways. Cubes about zero control
the cusp; small boxes away from the coordinate hyperplanes control
ordinary neighborhoods in the lattice space.

\begin{proposition}\label{prop:tightness}
Every sequence of the packet probabilities with $Q\to\infty$
is tight. More precisely, for $0<\theta<1$,
\begin{equation}\label{eq:cusp}
 \limsup_{Q\to\infty}\mu_*
 \{\Lambda:\exists\,0\ne v\in\Lambda,\ \|v\|_\infty<\theta\}
       \le C_n\theta^n.
\end{equation}
Every subsequential probability limit is $A_n$-invariant.
\end{proposition}

\begin{proof}
For $f=\one_{[-\theta,\theta]^n}$, the first $n-1$ logarithms
in \eqref{eq:product-fiber} are at most $\log\theta$.
The last coordinate bounds the sum of their gaps from
$\log\theta$ by $\log(\theta^n/z)$. Thus
\[
 \cV_f(z)=\frac{2^n}{(n-1)!}
           \bigl(\max\{0,\log(\theta^n/z)\}\bigr)^{n-1}.
\]
Put $X=\theta^nQ$ and $U=Q^{3/4}$. Up to an $n$-dependent
constant, Proposition~\ref{prop:unfolding} reduces the first
moment to
\[
 \frac1{\sqrt D\,\kappaK}
 \sum_{\N\mathfrak a\le X}
     P(\mathfrak a)\left(\log\frac X{\N\mathfrak a}\right)^{n-1}.
\]
If $X<1$, the sum is empty. Otherwise its logarithmic factor is
at most $(\log Q)^{n-1}$, since $\theta<1$. The terms with
$\N\mathfrak a\le\min(U,X)$ contribute at most
$C_nQ^{-1/4}(\log Q)^{n-1}$, by applying
Proposition~\ref{prop:weighted-count} at $x=y=U$.
When $X>U$, define the truncated cumulative sum
\[
 A_U(t)=\sum_{U<\N\mathfrak a\le t}P(\mathfrak a),
 \qquad U\le t\le X.
\]
Writing each logarithmic power as an integral and applying Tonelli gives
\begin{align*}
 \sum_{U<\N\mathfrak a\le X}
   P(\mathfrak a)\left(\log\frac X{\N\mathfrak a}\right)^{n-1}
 &= (n-1)\int_U^X A_U(t)
      \left(\log\frac Xt\right)^{n-2}\frac{dt}{t} \\
 &\ll_n \frac{\kappaK}{q}
      \int_U^X\left(\log\frac Xt\right)^{n-2}\dd t
 \ll_n \frac{\kappaK}{q}X.
\end{align*}
Here the same proposition at $x=y=t$ bounds $A_U(t)$, and the
last inequality follows from
$\int_0^1(\log(1/v))^{n-2}\dd v<\infty$.
There is no lower boundary term, since $A_U(U)=0$, and the
logarithmic weight vanishes at $X$. If $X\le U$, this part is absent.
Dividing by $\sqrt D\,\kappaK$ gives $O_n(\theta^n)$.
Since $E_f\ge1$ whenever the displayed cusp event occurs,
this proves \eqref{eq:cusp}.

Mahler's compactness criterion \cite[Section~3, Theorem~2]{Mahler1946} says that the
complement of this cusp event is compact. Given a desired
mass error, choose $\theta$ first and then $Q$ sufficiently large.
The finitely many earlier packet supports are compact by
Lemma~\ref{lem:packet-model}, so the whole sequence is tight.
The lattice space is locally compact and second countable;
weak compactness for tight probabilities gives probability
subsequential limits. Invariance passes to them by testing
against continuous compactly supported functions.
\end{proof}

For $r>0$, let
\[
 B(r)=\{h\in\SL_n(\R):
             \max_{i,j}|h_{ij}-\delta_{ij}|<r\}.
\]

\begin{proposition}\label{prop:ball-bound}
Let $\mu$ be a weak limit of packet probabilities along
$Q\to\infty$. For every compact $\Omega\subset X_n$ there are
$C_\Omega,r_\Omega>0$ such that
\begin{equation}\label{eq:ball-bound}
 \mu(xB(r))\le C_\Omega r^n
       \qquad(x\in\Omega,\ 0<r<r_\Omega).
\end{equation}
\end{proposition}

\begin{proof}
First fix a compact $V\subset(\R^\times)^n$ and $b_0\ge1$.
For $w\in V$ and small $r$, take
$f=\one_{w+[-b_0r,b_0r]^n}$. Each coordinate has a fixed sign,
is uniformly separated from zero, and has logarithmic width
$O_{V,b_0}(r)$. Its absolute product lies within $Cr$ of
$z_w=\prod_j|w_j|$, so
\[
 \cV_f(z)\le Cr^{n-1}
       \one_{[z_w-Cr,z_w+Cr]}(z)
\]
for some $C\ge1$ depending only on $V,b_0,n$.
Choose $r_0$ small enough that $z_w>4Cr_0$ on $V$.
Use Proposition~\ref{prop:weighted-count} with
\[
 x=(z_w+2Cr)Q,\qquad y=4CrQ.
\]
This slightly enlarged interval contains every contributing
ideal norm. Uniformly in $w$, for sufficiently large $Q$ it has
$Q^{3/4}\le x\le A_0Q$ with a fixed $A_0\ge1$, and
$\eta x\le y\le x$ with $\eta=4Cr/A_0\le1$, after decreasing
$r_0$ if necessary. For every fixed $r\in(0,r_0)$,
Proposition~\ref{prop:unfolding} therefore gives, eventually,
\begin{equation}\label{eq:vector-box}
 \int E_{\one_{w+[-b_0r,b_0r]^n}}\dd\mu_*
       \ll_{V,b_0,n}
       \frac{r^{n-1}}{\sqrt D\,\kappaK}
       \frac{\kappaK rQ}{q}=r^n .
\end{equation}
Here $Q=q\sqrt D$. The implicit constant is independent of $r$. The threshold for $Q$
may depend on $r$, which is all that will be needed.

We next choose a suitable lattice vector at each $x\in\Omega$.
Every full lattice has a vector with no zero coordinate:
for any basis, its combination with coefficients
$(1,N,\ldots,N^{n-1})$ has coordinates that are nonzero
polynomials in $N$, so an integer $N$ can avoid their finitely
many roots. Near each center choose a continuous local lift to basis
matrices, and continue the chosen vector using its fixed integer
basis coefficients. On a sufficiently small neighborhood with compact
closure in that chart, its coordinates remain bounded and bounded away
from zero. A finite cover of $\Omega$ by such neighborhoods therefore
gives choices $w_x$ lying in one compact
$V\subset(\R^\times)^n$. This pointwise selection need not be
continuous or measurable.

For $h$ in $B(r)$, each coordinate of $w_xh-w_x$ has absolute
value at most $r\sum_i|(w_x)_i|$. Thus one fixed $b_0$ makes
every lattice in $xB(r)$ contain a nonzero vector in
$w_x+[-b_0r,b_0r]^n$. Equation~\eqref{eq:vector-box} bounds
the packet mass of $xB(r)$ by $C_\Omega r^n$.
This set is open, so Portmanteau gives
\[
 \mu(xB(r))\le\liminf_i\mu_i(xB(r))\le C_\Omega r^n.
\]
The argument holds for every fixed $x$ and $r$ with the same
constant on $\Omega$. It proves \eqref{eq:ball-bound} for
all $0<r<r_\Omega$, without any exchange of shrinking-radius
and packet limits.
\end{proof}

\section{Entropy and identification of the limit}\label{sec:entropy}

The ball estimate has exponent $n$, whereas $A_n$ has dimension $n-1$.
We follow the ordinary-ball strategy of
\cite[Section~2.7.1]{ELMV2011}. We retain the estimate for every invariant
probability dominated by a finite multiple of the limit, so it forces
positive entropy after restriction to any positive-weight collection of
ergodic components. Measure classification then identifies the limit.

For an invariant probability $\nu$, write $h_\nu(a)$ for the
measure-theoretic entropy of the map $x\mapsto xa$.
Fix distinct real numbers $\lambda_1,\ldots,\lambda_n$ with sum zero and put
\[
 a(t)=\operatorname{diag}(e^{\lambda_1t},\ldots,e^{\lambda_nt}),
 \qquad \lambda=\min_{i\ne j}|\lambda_i-\lambda_j|>0.
\]

\begin{proposition}[Entropy from ordinary balls]\label{prop:entropy-from-balls}
Let $n\ge3$ and let $\mu$ be an $A_n$-invariant probability on $X_n$.
Suppose that for every compact $\Omega\subset X_n$ there are
$C_\Omega,r_\Omega>0$ such that
\[
 \mu(xB(r))\le C_\Omega r^n
 \qquad(x\in\Omega,\ 0<r<r_\Omega).
\]
Then every $A_n$-invariant probability $\xi$ satisfying $\xi\le c\mu$
for some finite $c$ has $h_\xi(a(1))\ge\lambda/3$.
\end{proposition}

\begin{proof}
Fix $0<b<1/4$ and consider the two-sided neighborhoods
\[
 T_t=a(t)B(b)a(-t)\cap a(-t)B(b)a(t),\qquad t\ge0.
\]
If $h\in T_t$, then
\begin{equation}\label{eq:tube-entries}
 |h_{ii}-1|<b,\qquad
 |h_{ij}|<b e^{-|\lambda_i-\lambda_j|t}\quad(i\ne j).
\end{equation}
Set $r=e^{-\lambda t}$.  Every term in the determinant expansion other
than the product of the diagonal entries contains at least two
off-diagonal entries.  Since $\det h=1$, it follows that
$\prod_i h_{ii}=1+O_{n,b}(r^2)$.  Thus
\[
 d(h)=\operatorname{diag}\left(h_{11},\ldots,h_{n-1,n-1},
                   \Bigl(\prod_{i<n}h_{ii}\Bigr)^{-1}\right)
\]
belongs to a fixed compact subset $\mathcal D\subset A_n$, and
$d(h)^{-1}h\in B(Cr)$ for a constant depending only on $n,b$.

In logarithmic coordinates, $\mathcal D$ is a bounded subset of $\R^{n-1}$.
Covering it by a mesh of size comparable to $r$, and enlarging $C$ to
allow for matrix multiplication, gives
\begin{equation}\label{eq:tube-cover}
 T_t\subset\bigcup_{j=1}^{N_t}d_jB(Cr),
 \qquad d_j\in\mathcal D,\qquad N_t\le Cr^{-(n-1)}.
\end{equation}
For $x$ in a compact set $\Omega$, all the centers $xd_j$ lie in the
fixed compact set $\Omega\mathcal D$.  The assumed ball estimate and domination
therefore imply, for all sufficiently large $t$,
\begin{equation}\label{eq:tube-mass}
 \xi(xT_t)\le C_{\Omega,b,c}\,r^{-(n-1)}r^n
             =C_{\Omega,b,c}e^{-\lambda t}
 \qquad(x\in\Omega).
\end{equation}
Since $\xi$ is a probability, the constant can be enlarged to cover
the finitely many integer times before this estimate holds.

We use the entropy criterion of
Einsiedler--Lindenstrauss--Michel--Venkatesh
\cite[Corollary~3.3]{ELMV2009}.  In the notation
\[
 B^{(s,t)}=a(-s)Ba(s)\cap a(-t)Ba(t),
\]
the criterion states the following: if flow-invariant probabilities
$\xi_i$ converge weakly to a probability $\xi$, and $t_i\to\infty$,
then $h_\xi(a(1))\ge\eta$ provided that, for every compact set $\Omega$,
there are an identity neighborhood $B$ and a constant $C_\Omega$ such that
\[
 \xi_i\bigl(xB^{(-t_i,t_i)}\bigr)
       \le C_\Omega e^{-2\eta t_i}\qquad(x\in\Omega).
\]
This criterion applies on the noncompact lattice space; its requirement
that the limit be a probability is essential.  Here take the constant
sequence $\xi_i=\xi$, the neighborhood $B(b)$, $t_i=i$, and
$\eta=\lambda/3$.  Since $B(b)^{(-t,t)}=T_t$,
Equation~\eqref{eq:tube-mass} proves the required estimate.
\end{proof}

\begin{lemma}[Entropy in the ergodic components]\label{lem:entropy-components}
Under the hypotheses of Proposition~\ref{prop:entropy-from-balls},
almost every $A_n$-ergodic component of $\mu$ has positive entropy
for $a(1)$.
\end{lemma}

\begin{proof}
The group $A_n\simeq\R^{n-1}$ is locally compact and second countable,
and its action on the standard Borel space $X_n$ is continuous.
The ergodic decomposition theorem
\cite[Theorem~5.2]{GreschonigSchmidt2000}, specialized to invariant
probabilities, therefore supplies a measurable decomposition
\[
 \mu=\int_Z\nu_z\,\dd\tau(z)
\]
into $A_n$-invariant, $A_n$-ergodic probabilities over a factor fixed
by $A_n$. Set $T(x)=xa(1)$.
The function $z\mapsto h_{\nu_z}(T)$ is measurable.  Indeed, choose
increasing finite Borel partitions $\mathcal P_k$ generating the Borel
$\sigma$-algebra of $X_n$.  For every $T$-invariant probability $\nu$,
\[
 h_\nu(T)=\sup_k\inf_{N\ge1}\frac1N
 H_\nu\left(\bigvee_{j=0}^{N-1}T^{-j}\mathcal P_k\right),
\]
where $H_\nu$ denotes the Shannon entropy of a finite partition.
The right side is a countable combination of Borel functions of $\nu$.

The component partition is fixed by $T$, since $a(1)\in A_n$.
The entropy integral formula for a fixed invariant partition
\cite[Theorem~9.8]{Rokhlin1967} consequently gives
\begin{equation}\label{eq:entropy-decomposition}
 h_\mu(T)=\int_Z h_{\nu_z}(T)\,\dd\tau(z).
\end{equation}
The same formula holds after restricting to any measurable collection
of components.  It does not require the $\nu_z$ to be $T$-ergodic:
the induced transformation on the component space is the identity
and contributes zero entropy.

If $Z_0=\{z:h_{\nu_z}(T)=0\}$ had weight $w=\tau(Z_0)>0$, then
\[
 \xi=\frac1w\int_{Z_0}\nu_z\,\dd\tau(z)
\]
would be an $A_n$-invariant probability with $\xi\le w^{-1}\mu$.
Equation~\eqref{eq:entropy-decomposition}, applied to this restricted
decomposition, would give $h_\xi(T)=0$, contradicting
Proposition~\ref{prop:entropy-from-balls}.
\end{proof}

\paragraph{Measure classification.}
Einsiedler--Katok--Lindenstrauss
\cite[Theorem~1.3 and Corollary~1.4]{EKL2006} prove that, for prime
$n\ge3$, an $A_n$-invariant and $A_n$-ergodic probability on
$\operatorname{SL}_n(\R)/\operatorname{SL}_n(\Z)$ is Haar whenever
some diagonal one-parameter subgroup has positive entropy.
Their action is on the left.  The map
$\Gamma g\mapsto g^{-1}\Gamma$, where $\Gamma=\operatorname{SL}_n(\Z)$,
changes our right action by $a(t)$ into the left action by $a(-t)$.
Since entropy is unchanged by inversion of an invertible transformation,
their result applies to our convention without any additional hypothesis.

\begin{proof}[Proof of Theorem~\ref{thm:main}]
Let $\mu_i$ be the packet probabilities for a sequence in the theorem.
Equation~\eqref{eq:discriminant-scale} gives $Q_i\to\infty$.
By Proposition~\ref{prop:tightness}, this sequence is tight and every
weak subsequential limit $\mu$ is an $A_n$-invariant probability.
Proposition~\ref{prop:ball-bound} supplies the hypotheses of
Proposition~\ref{prop:entropy-from-balls}.
Lemma~\ref{lem:entropy-components} shows that almost every
$A_n$-ergodic component of $\mu$ has positive entropy for $a(1)$.
As $n$ is prime, the classification just stated makes every such
component equal to $m_n$.  Integrating the components gives $\mu=m_n$.

Thus every convergent subsequence has the same probability limit.
Together with tightness, this proves $\mu_i\to m_n$ and excludes
loss of mass.  All preceding estimates were for each prescribed local
homothety packet separately, so no averaging over different local
module types is used.
\end{proof}


\appendix
\section{Adelic formulation and choices}\label{subsec:adelic-packet}

The model of Lemma~\ref{lem:packet-model} also identifies the measure obtained from adelic
torus Haar probability. We give the identification for arbitrary local
type, so the lattice formulation includes all finite lattice and real
diagonalization choices in the usual adelic formulation.

Let $\mathbb A_f$ be the finite adeles of $\Q$, let
$\widehat\Z=\prod_\ell\Z_\ell$, and put
$\mathbb A_\Q=\R\times\mathbb A_f$ and
$\mathbb A_K=K\otimes_\Q\mathbb A_\Q$.
Write $\widehat R=R\otimes\widehat\Z$.
Choose a rational basis $b_1,\ldots,b_n$ of $K$, and let
$\rho(\alpha)$ denote multiplication by $\alpha$ on row coefficient
vectors. It embeds the torus
\[
 \mathbf T_K=(\operatorname{Res}_{K/\Q}\mathbb G_m)/\mathbb G_m
 \quad\hbox{in}\quad G=\operatorname{PGL}_{n/\Q}.
\]
The embedding matrix $B_\sigma=(\sigma_j(b_i))_{i,j}$ satisfies
\begin{equation}\label{eq:regular-diagonalization}
 \rho(\alpha)B_\sigma
   =B_\sigma\operatorname{diag}(\sigma_1(\alpha),\ldots,\sigma_n(\alpha)).
\end{equation}
For $g_f\in G(\mathbb A_f)$, choose a lift to
$\operatorname{GL}_n(\mathbb A_f)$ and set
\[
 \widehat M(g_f)=\widehat\Z^{\,n}g_f^{-1},\qquad
 M(g_f)=K\cap\widehat M(g_f),
\]
using the chosen rational basis to identify coefficient vectors with
elements of $K$. This is a full lattice, and its completion is
$\widehat M(g_f)$: choose an integer $d\ge1$ with
$d\widehat\Z^{\,n}\subset\widehat M(g_f)
\subset d^{-1}\widehat\Z^{\,n}$, and patch in the finite quotient
$d^{-1}\Z^n/d\Z^n$. Changing the projective lift multiplies it by a
scalar finite idele. Every such idele is a rational scalar times an
element of $\widehat\Z^\times$; the latter preserves every completed
full lattice. Consequently this change only rescales $M(g_f)$ by
a rational number.

Choose $g_\infty\in G(\R)$ with
$g_\infty^{-1}\mathbf T_K(\R)^\circ g_\infty=A_n$, using the
odd-degree identification of $G(\R)$ with $\SL_n(\R)$.
For $g=(g_\infty,g_f)$, the adelic packet is the projection of
$\mathbf T_K(\Q)\backslash\mathbf T_K(\mathbb A_\Q)g$ to
$G(\Q)\backslash G(\mathbb A_\Q)/G(\widehat\Z)$, equipped with
the pushforward of torus Haar probability.

\begin{proposition}\label{prop:adelic-packet}
The adelic torus quotient is compact. Its packet and canonical
probability agree with $\mathcal P_{K,M(g_f),\sigma'}$ and
$\mu_{K,M(g_f),\sigma'}$ for an ordering $\sigma'$ of the real
embeddings. Every prescribed local homothety type occurs this way.
Changing the rational basis or the projective lifts does not change
the packet; changing the allowed real diagonalizer only reorders the
embeddings. Thus Theorem~\ref{thm:main} includes every such adelic choice.
\end{proposition}

\begin{proof}
Write $M=M(g_f)$ and $\widehat M=\widehat M(g_f)$.
First identify the ambient quotient with $X_n$. Given a finite matrix
lift $g_f$, choose a basis matrix $c\in\operatorname{GL}_n(\Q)$
for $\Q^n\cap\widehat\Z^{\,n}g_f^{-1}$. Then
$cg_f\in\operatorname{GL}_n(\widehat\Z)$, and the corresponding
real lattice is $\Z^ncg_\infty$, normalized to covolume one.
Changing this basis or either matrix lift has no effect on that
lattice. This constructs the usual quotient identification directly.
For odd $n$, the signed real $n$th root gives
$[h]\mapsto(\det h)^{-1/n}h$, identifying
$G(\R)$ with $\SL_n(\R)$ and $G(\Z)$ with $\SL_n(\Z)$.

Multiplication matrices modulo scalars give
\begin{equation}\label{eq:idelic-torus-quotient}
 \mathbf T_K(\Q)\backslash\mathbf T_K(\mathbb A_\Q)
 \simeq K^\times\backslash\mathbb A_K^\times/\mathbb A_\Q^\times.
\end{equation}
Here the pointwise lifts can also be seen by elementary linear algebra:
over any ground field $F_0$ containing $\Q$, a matrix whose projective
class lies in $\mathbf T_K$ becomes a scalar multiple of a multiplication
operator over an algebraic closure. It therefore commutes exactly with
the regular representation already over $F_0$. The centralizer of that
representation is the algebra of multiplication operators, so the matrix
belongs to $(K\otimes_\Q F_0)^\times$, with scalar ambiguity $F_0^\times$.
Finite lifts may be chosen integral at almost all primes.

A finite idele $t_f$ acts on the completed lattice by
$\widehat M\mapsto t_f^{-1}\widehat M$. Hence its translates are
exactly the prescribed local homothety type. The stabilizer of
$\widehat M$ among finite ideles is
\[
 \{z:z\widehat M=\widehat M\}
       =(\cO(M)\otimes\widehat\Z)^\times:
\]
equality holds precisely when both $z$ and $z^{-1}$ preserve
$\widehat M$. In the projective torus the stabilizer is the image
of this group, since a representative that preserves the lattice
up to a scalar idele can be rescaled to preserve it exactly.
The unit subgroup is compact open in the finite ideles and has finite
index in $\widehat R^\times$; at each prime not dividing $q$ it equals
$R_\ell^\times$. Its projective image is compact open as well.
No invertibility of $M$ over $\cO(M)$ is used.

To describe the quotient and its measure, first divide
\eqref{eq:idelic-torus-quotient} by the image of
$\widehat R^\times$. Fractional-ideal factorization and
$\mathbb A_\Q^\times
=\Q^\times\R^\times\widehat\Z^\times$
give one real fiber $E\backslash C$ for each ordinary ideal class.
These fibers are compact by the unit theorem, and the omitted
finite-unit group is compact. Thus the torus quotient is compact.
At the stabilizer of $\widehat M$, the finite fibers instead retain
the translates $\mathcal N(I)$. The same decomposition becomes
\[
 \coprod_{[I]}(\mathcal N(I)\times C)/E,
\]
with exactly the action in Lemma~\ref{lem:packet-model}.
Haar measure pushes to uniform counting measure on each finite
unit orbit and to the stated logarithmic Haar measure on $C$.
After normalization it is therefore precisely the probability in
that lemma. For $g_\infty=[B_\sigma]$,
\eqref{eq:regular-diagonalization} identifies the resulting real
lattices with its displayed map. This proves both the packet
identification and the equality of measures, including the orbit
volume weights.

Finally, every full $M$ is obtained by choosing a basis matrix for
its completed lattice and taking its inverse as $g_f$.
Any other allowed real diagonalizer has the form $[B_\sigma dP]$,
where $d$ is diagonal and $P$ is a permutation matrix, because it
must permute the joint eigenspaces. The diagonal factor translates
the real torus and permutes its sign components, preserving its
Haar probability; the permutation reorders the embeddings.
If a rational change of basis has matrix $c$, then
$B'_\sigma=cB_\sigma$ and $\rho'=c\rho c^{-1}$, while the same
packet data have representatives $g'_\infty=cg_\infty$ and
$g'_f=cg_f$. Thus the torus is conjugated and the translated packet
is left-multiplied by $c$, which disappears in the ambient rational quotient.
\end{proof}


{\small
\bibliographystyle{plain}
\bibliography{references}
}
\end{document}
